\documentclass[10pt]{article}

\usepackage[utf8]{inputenc}

\usepackage[T1]{fontenc}

\usepackage{graphicx}

\usepackage[export]{adjustbox}

\graphicspath{ {./images/} }

\usepackage{amsmath}

\usepackage{amsfonts}

\usepackage{amssymb}

\usepackage[version=4]{mhchem}

\usepackage{stmaryrd}



\begin{document}

встречающихся вачально-краевых задач (см. Краевая задача; численные методы решения для уравнений с частными производными), приближенного вычисления интегралов (см. Интегрирование численное). В результате исходная задача оптимального управления заменяется нек-рым семейством аппроксимирующих задач, зависящим от нек-рых параметров (напр., от пагов разностной сетки). Вопросы построения аппроксимирующих задач, исследование сходимости см. в [5].



Широкие классы әкстремальных задач являются некорректно поставленными (см. Некорректные задачи) и для их решения нужно использовать регуляризации методы (см. [5], [13]).



Лum.: [1] А ле е к с е e в B. М., Т их х ом и р о в B. М.,



\includegraphics[max width=\textwidth, center]{2023_07_09_0b148d5e28a13f3f5e3dg-1}

алгоритмы решения задач оптимивации, К., 1983; [3] Б у т к о вс к и й А. Г., Методы управления системами с распределенными параметрами, М., 1975; [4] В а с и л в е в Ф. П.2 Численные методы решения әкстремальных задач, М., 1980; [5] е г о ж е, Методы решения әкстремальных задач, М., 1981; [6] Е в т у ша е нк о Ю. Г., Методы решения әкстремальных задач и их применение в системах оптимизации, М., 1982; [7] Е г о р о в А. И., Оптимальное управление тепловыми и дифффузионными процессами, М., 1978; [8] К р а с о в с к и й Н. Н., Теория управления движением, М., 1968; [9] ЈI и о \& с Ж.-ЈI., Оптимальное управление системами, описываемыми уравнениями с частными производными, пер. с франц., М., 1972; [10] П о л я к Б. Т., Введение в оптимизацию, М., 1983; [11] С е а Ж., Оптимизация. Теория и алгоритмы, пер. с франц., М., 1973; [12] С и р а з е тд и н о в Т. $\mathrm{K}$. Оптимизация систем с распределенными параметрами, М., 1977; [13] Т и х о н о в А. Н., А р с е н и н В. Я. Методы решения некорректных задач, 2 изд., М., 1979; [14] $\Phi$ е д о р е н к о Р. П., Приближенное решение задач оптимального управления, М., 1978; [15] Э к л а є д И., Т е м а м P., М., 1979.



ЭКСТРЕМАЛЬНЫЕ СВОЙСТВА ПОЛИНОМОВ свойства алгебраических, тригонометритеских или обобщенных полиномов, к-рые выделяют их в качестве решөний нек-рых әкстремальных задач.



Напр., Чебышева многочлены $T_{n}(x)=\cos (n \arccos x)=$ $=2^{n-1} x^{n}+\ldots$ имеют наименьшую норму в пространстве $C([-1,1])$ среди всех алгебраич. многочленов степени $n$ со старшим коэффициентом, равным $2^{n-1}$ (П. Л. Чебышев, 1853) и, таким образом, являются решением следующей экстремальной задачи



\[

\max _{x \in[-1,1]}\left|2^{n-1} x^{n}+a_{1} x^{n-1}+\ldots+a_{n}\right| \longrightarrow_{a=\left(a_{1}, \ldots, a_{n}\right)} \inf ^{1}

\]



Иначе говоря, многочлен $T_{n}$ наименее уклоняется от нуля в пространстве $C([-1,1])$ среди всех многочленов степени $n$ со старшим коэффициентом, равным $2^{n-1}$.



Экстремальные задачи на пространствах многочлевов в значительной части исследуются в пространствах $L_{p}([a, b]), 1 \leqslant p \leqslant \infty$. При этом наиболее известные результаты связаны со случаями $p=1,2$ и $\infty$ (метрика $C$ ). В частности, в этих метриках решен вопрос о явном виде многочленов, наименее уклоняощихся от нуля. В метрике $L_{1}$ - это многочлены Чебышева 2-го рода, в метрике $L_{2}$ - Лежаидра жиогочленьц, о метрике $C$ см. выше. Описано также множество классических ортогональных многочленов, являющихся многочленами наименее уклоняющимися от нуля в пространстве $L_{2}$ с весом (Лагерра жногочлены, Эрлита многочлены, Якоби многочленьи и т. ПІ.).



Е. И. Золотарев (1877) рассмотрел вонрос о многочленах вида $x^{n}+\sigma x^{n-1}+a_{2} x^{n-2}+\ldots+a_{n}$ (с двумя фиксированными старшими коэфффициентами), наименее уклоняющихся от нуля в метрике $C$. Он нашел однопараметрич. семейство многочленов, решающих әту задачу, выразив их через эллиптич. функции.



Многочлены Чебышева экстремальны в задаче о неравенстве для производных, а именно, имеет место следующее точное Маркова неравенство (где $P_{n}-$ многочјен степени $\leqslant n)$ :



\[

\left\|P_{n}^{(k)}\right\|_{C([-1,1])} \leqslant\left|T_{n}^{k}(1)\right|\left\|P_{n}\right\|_{C([-1,1])},

\]



в к-ром равенство достигается на $T_{n}$. Неравенство (*) для $k=1$ доказал А. А. Марков (1889), для остальных $k$ - B. А. Марков (1892). О подобном неравенстве для тригонометрич. полиномов см. Бернщтейна неравенcmeo.



Нек-рые әкстремальные свойства алгебраических и тригонометрических полиномов в равномерной метрике переносятся на чебышевские системы функций (см. [2]). О теории экстремальных задач и Ә. с. п. см. [6].



M. JIит.: [1] ч е 6 ы ш е в тг. Ј., Полын. собр. соч., т. 2-3, свойства полиномов и наилучшее приближение негрерывных функций одной вещественной переменной, ч. 1, Л.- М., 1937; [3] Го н ч а р о в В. Л., Теория интерполирования и приблияения функций, 2 изд., М., $1954 ;[4]$ А х и е з е р Н. И., Јекции по теорпи аппроксимации, 2 изд., М., 1965 ; [5] В о р о-, н о в с к а ғ Е. В., Метод фугкционалов и его приложения ЈI., 1963; [6] Т и х о м и р о в В. М., Некоторые вопросы тео-



ЭКСТРЕМАЛЬНЫЕ СВОЙСТВА ФУ НКЦИЙ - своЙства отдельных функций, к-рые выделяют их как решения нек-рых әкстремальных задач. Большинство специальных функций, возникших в математич. анализе могут быть охарактеризованы нек-рым экстремальным свойством. Таковы, напр., әкстремальные свойства полиномов: классич. Лагерра многочлены, Лежандра многочлены, Чебышева многочлены, Эрмита многочлены, Якоби многочлень можно охарактеризовать как многочлены, наименее уклоняющиеся от нуля в пространстве $L_{2}$ с весом. Классич. полиномы являются обычно решениями разных экстремальных задач, нередко возникающих в отдаленных областях анализа. Так, напр., многочлены Чебышева әкстремальны в задаче о неравенстве для производных многочленов (см. [1], Маркова неравенство). То же можно сказать и о др. специальных функциях. Многие из них являются собственными функциями для дифференциальных операторов, т. е. являются репением нек-рой изопериметрической задачи. При этом, наиболее известные специальные функции так или иначе связаны с наличием нек-рой инвариантной структуры (см. Гармонический анализ абстрактньй), когда они являются собственными функциями JI апласа - Бельтрами уравнения, инвариантного относительно сдвигов. Таковы тригонометрич. полиномы, сферич. Функции, цилиндрич. функции и др. (см. [2]). Большинство Э. с. ф. может быть сформулировано в виде нек-рого точного неравенства.



C экстремальными задачами теории приближений связаны Бернштейна неравенство, Бора - Фаваранеравенстөо и др. В частности, неравенство Бора - Фавара отражает экстремальное свойство Бернулли лногочленов. Э. с. Ф. изучаются в теории приближений (см. [6], [7]), в теории численного интегрирования (см. [8]). Сплайны могут быть охарактеризованы различными экстремальными свойствами (см. [9]). Многие специальные сплайны обладают рядом экстремальных свойств, касающихся ашпроксимации и интерполяции классов функций (см. [7], [8]). Многие Ә. с. ф. изучают в комплексном анализе. В частности, Кёбе функция является экстремальной функцией рида задач теории однолистных функций. См. также Изопериметрическое неравенство, Вломения теоремы.



Jит.: [1] Б е р н ш т е й н С. Н., Экстремальные свойства полиномов и наилучшее приближение непрерывных функций одной вещественной переменной, ч. 1, ЈI.- М., 1937; [2] В иле ен к и н Н. Я., Специальные функции и теория представлений трупII, М., 1965; [3] X а р д и $\Gamma$. Г., ЈІ ит т Л bв у д Д. Е., ІІ о ли а Г., Неравенства, иер. с англ., М., 1948 [4] Бе к к е н б а х Э., 'Бе л л м а н Р., Неравенства, пер. c англ., M., 1965; [5] M i t r i n ov i č D. S., Analitičke nejednakosti, Beograd, 1970; [6] К о р н е й ч у к H. П., Экстремальные задачи теории приближений, М., 1976; [7] Т й х о м иро в B. М., Некоторые вопросы теории приближений, М., 1976 ; [8] Н и 1 0 л b с к и й С. М., Квадратурные формулы, 3 изд., М., 1979; [9] А л 6 е р г Д ж., Н и л ь с о н Э., У о л Д. ж., Теория сплайнов и его приложения, пер. с англ., М., 1972



B. $M$. Тихолиров.





\end{document}